\documentclass{beamer}

\usepackage[size=a0,scale=1.15]{beamerposter}

% Theme (safe choice)
\usetheme{Simple}

% Encoding + fonts
\usepackage[utf8]{inputenc}
\usepackage[T1]{fontenc}
\usepackage{lmodern}

% Graphics
\usepackage{graphicx}
\usepackage{booktabs}

% Remove navigation symbols (fixes .nav issues)
\setbeamertemplate{navigation symbols}{}

% Title info
\title{False Memories to Fake News:\\
The Evolution of Misinformation in Academic Literature}
\author{Ruiz Iglesias et al.}
\institute{University of Vermont \\ Computational Story Lab}

\begin{document}

\begin{frame}[t]

\begin{columns}[t,totalwidth=\linewidth]

% -------- LEFT --------
\begin{column}{0.32\linewidth}

\begin{block}{Big Idea}
\begin{itemize}
\item Misinformation research surged after 2016
\item Rooted in earlier psychology work
\item Links memory studies to social media
\end{itemize}
\end{block}

\begin{block}{Research Questions}
\begin{itemize}
\item How did the term evolve?
\item What traditions shaped it?
\item What is missing today?
\end{itemize}
\end{block}

\begin{block}{Data \& Methods}
\begin{itemize}
\item Papers (2011--2023)
\item Keyword analysis
\item Citation networks
\end{itemize}
\end{block}

\end{column}

% -------- CENTER --------
\begin{column}{0.36\linewidth}

\begin{block}{Key Results}

\textbf{Explosion After 2016}
\begin{itemize}
\item Rapid growth of publications
\end{itemize}

\vspace{0.5em}

\textbf{Two Paradigms}
\begin{itemize}
\item Modern: social media spread
\item Older: cognition and memory
\end{itemize}

\vspace{0.5em}

\textbf{Key Insight}
\begin{itemize}
\item These traditions are not fully integrated
\end{itemize}

\end{block}

\begin{block}{Figures}
\centering
\includegraphics[width=0.7\linewidth]{example-image-a}

\vspace{1em}

\includegraphics[width=0.7\linewidth]{example-image-b}
\end{block}

\end{column}

% -------- RIGHT --------
\begin{column}{0.32\linewidth}

\begin{block}{Concept}
\begin{itemize}
\item A “misinformation paradigm”
\item Combines cognition + networks
\end{itemize}
\end{block}

\begin{block}{Implications}
\begin{itemize}
\item Bridge disciplines
\item Rethink assumptions
\end{itemize}
\end{block}

\begin{block}{Conclusion}
\begin{itemize}
\item Field is newer than it seems
\item Built on older psychology research
\end{itemize}
\end{block}

\begin{block}{Takeaway}
\centering
\large
Misinformation research blends psychology and social systems
\end{block}

\end{column}

\end{columns}

\end{frame}

\end{document}\documentclass[final]{beamer}

\usepackage[size=a0,scale=1.2]{beamerposter}
\usetheme{default}

\usepackage{graphicx}
\usepackage{amsmath}
\usepackage{booktabs}

\title{False Memories to Fake News:\\
The Evolution of “Misinformation” in Academic Literature}
\author{Ruiz Iglesias et al.}
\institute{University of Vermont \\ Computational Story Lab}

\begin{document}

\begin{frame}[t]

\begin{columns}[t]

% ---------------- LEFT COLUMN ----------------
\begin{column}{0.32\linewidth}

\block{Big Idea}{
\begin{itemize}
\item “Misinformation” is not just a modern concept
\item Research surged after 2016
\item Roots in psychology and memory research
\item Connects false memory studies to social media dynamics
\end{itemize}
}

\block{Research Questions}{
\begin{itemize}
\item How has “misinformation” evolved?
\item What traditions shaped it?
\item Are current narratives incomplete?
\end{itemize}
}

\block{Data \& Methods}{
\textbf{Dataset:}
\begin{itemize}
\item Academic papers (2011–2023)
\item Keywords: misinformation, disinformation, fake news
\end{itemize}

\textbf{Methods:}
\begin{itemize}
\item Term frequency analysis
\item Citation networks
\item Community detection
\end{itemize}
}

\end{column}

% ---------------- CENTER COLUMN ----------------
\begin{column}{0.36\linewidth}

\block{Key Results}{

\textbf{Explosion After 2016}
\begin{itemize}
\item Rapid growth in publications
\item Linked to social media, politics, pandemics
\end{itemize}

\vspace{1em}

\textbf{Two Research Traditions}

\textbf{Modern Paradigm}
\begin{itemize}
\item Focus: large-scale spread
\item Social networks and platforms
\end{itemize}

\textbf{Older Paradigm}
\begin{itemize}
\item Focus: cognition and memory
\item False recall, suggestibility
\end{itemize}

\vspace{1em}

\textbf{Hidden Continuity}
\begin{itemize}
\item Traditions developed separately
\item Still not fully integrated
\end{itemize}
}

\block{Figures}{
\centering
% Replace with actual figures
\includegraphics[width=0.8\linewidth]{example-image}
\\
\small Time series of misinformation research

\vspace{1em}

\includegraphics[width=0.8\linewidth]{example-image}
\\
\small Network of research communities
}

\end{column}

% ---------------- RIGHT COLUMN ----------------
\begin{column}{0.32\linewidth}

\block{Conceptual Contribution}{
\begin{itemize}
\item Defines a “misinformation paradigm”
\item Combines:
\begin{itemize}
\item Individual cognition
\item Social information systems
\end{itemize}
\item Not a neutral concept—shaped by assumptions
\end{itemize}
}

\block{Implications}{

\textbf{For Research}
\begin{itemize}
\item Bridge cognitive and social approaches
\end{itemize}

\textbf{For Society}
\begin{itemize}
\item Avoid over-focusing on social media
\end{itemize}

\textbf{For Theory}
\begin{itemize}
\item Treat misinformation as a constructed lens
\end{itemize}
}

\block{Conclusion}{
\begin{itemize}
\item Field expanded rapidly after 2016
\item Deep roots in psychology
\item Modern work builds on earlier paradigms
\item Future: integrate cognition + networks
\end{itemize}
}

\block{Takeaway}{
\centering
\Large
“Misinformation research is a fusion of cognitive psychology and social media science.”
}

\end{column}

\end{columns}

\end{frame}

\end{document}
